% GENERATED, edit the French unit or the translation table
\mxcorrections{3}\label{C03-CORR}
\textbf{Teacher use.} The answers below concern the tasks not worked in the lesson. First have students state their method; give graded help before the solution. The durations are estimates to be calibrated in class.
\subsection*{Diagnostic and remediation}
\begin{enumerate}
\item $\pt{3;-4}$. \textit{Action:} Subtract end minus start on two drawings.
\item $\sqrt{3^2+(-4)^2}=5$. \textit{Action:} Review Pythagoras and distinguish norm from coordinates.
\item $\vect{AC}$. \textit{Action:} Draw the composite displacement.
\item $a^2+2ab+b^2$. \textit{Action:} Identify the cross term.
\item $\dfrac{3}{5}$. \textit{Action:} Name the sides relative to the angle.
\item $\pt{1;2}$. \textit{Action:} Average the two coordinates separately.
\end{enumerate}
\subsection*{Investigation activity}
Take points on the circle with diameter AB; proof via the midpoint I: product $MI^2-9$. The endpoints are included because one of the vectors is zero.
\subsection*{Immediate applications}
\begin{enumerate}
\item $0$ then $-10$. The cosine is 0 and $-1$ respectively.
\item $3t-12=0$, so $t=4$.
\item The square is $4-2+9=11$; the norm is $\sqrt{11}$.
\item $32=16+16-32\cos A$, so $\cos A=0$ and $A=90^\circ$.
\item $MI^2-9=7$: circle with centre $\pt{3;0}$ and radius 4.
\end{enumerate}
\subsection*{Exercises: method and results}
\paragraph{Exercise 3.1 -- Direct calculations.} Product $4-6=-2$; norms $\sqrt5$ and 5; sum $\pt{5;1}$ with square 26. Check: $5+2(-2)+25=26$.
\paragraph{Exercise 3.2 -- A right triangle?} $\vect{AB}=\pt{3;1}$, $\vect{AC}=\pt{-1;3}$; product 0. $AB=AC=\sqrt{10}$, area $\dfrac{AB\times AC}{2}=5$ area units.
\paragraph{Exercise 3.3 -- Length without a frame.} $\vect{AC}=\vect{AB}+\vect{AD}$ and $\vect{BD}=\vect{AD}-\vect{AB}$. Product $\dfrac{4\times 3}{2}=6$; $AC^2=16+9+12=37$ and $BD^2=16+9-12=13$: the diagonal $[AC]$ is the longest, which corresponds to an acute angle $\widehat{BAD}$ (positive dot product).
\paragraph{Exercise 3.4 -- Recovering an angle.} $\cos \theta =-\dfrac{6}{12}=-\dfrac{1}{2}$; in $\intff{0}{180^\circ}$, the angle is $120^\circ$. The domain ensures uniqueness.
\paragraph{Exercise 3.5 -- A variable locus.} $MI^2=k+\dfrac{AB^2}{4}$. If the right-hand side is negative: empty set; zero: the single point I; positive: circle whose radius is its square root. The discussion bears on a squared distance.
\paragraph{Exercise 3.6 -- Work done by a force.} $W=\dfrac{50\times 4\times 1}{2}=100$ J. At $120^\circ$, $W=-100$ J. The sign distinguishes help to the displacement from opposition; it does not describe a negative distance.
\paragraph{Exercise 3.7 -- Contrapositive.} The contrapositive is ``if $\vec u$ and $\vec v$ are orthogonal, then $\norm{\vec u+\vec v}^2=\norm{\vec u}^2+\norm{\vec v}^2$''; it follows from the expansion, since $\norm{\vec u+\vec v}^2=\norm{\vec u}^2+2\vec u\cdot\vec v+\norm{\vec v}^2$ and $\vec u\cdot\vec v=0$. It is equivalent to the original statement. We can also write $2\vec u\cdot\vec v\ne0$.
\paragraph{Exercise 3.8 -- Delayed review.} Al-Kashi: $\dfrac{9+16-25}{2}=0$. Converse of Pythagoras: triangle right-angled at A, so zero product. Name the hypotheses of each method.
\paragraph{CQFD, exercise 9 p.~250, questions a and b.} a: $1\times3+2\times(-1)=1$. b: $(-1)\times12+3\times7=9$. The scalar product is a number, not a vector.
\paragraph{CQFD, exercise 11 p.~250, questions 1, 2 and 3.} $\vect{AB}=\pt{4;6}$, $\vect{AC}=\pt{-3;4}$, $I\pt{-1;-2}$ and $\vect{CI}=\pt{5;-1}$. Hence $\vect{AB}\cdot\vect{AC}=12$, $\vect{AB}\cdot\vect{CI}=14$. Chasles gives $\vect{AB}\cdot\vect{AI}=12+14=26$. Two checks: $\vect{AI}=\pt{2;3}$ gives $8+18=26$; or $\vect{AI}=\dfrac12\vect{AB}$ gives $\dfrac12(16+36)=26$.
\subsection*{MCQ: answers and reasons}
\textbf{MCQ 1.}
\begin{enumerate}
\item A. $6-4=2$.
\item C. Square root of 25.
\item D. Scalar criterion.
\item B. $\dfrac{2\times 3}{2}=3$.
\item C. Expand the scalar square.
\end{enumerate}
\textbf{MCQ 2.}
\begin{enumerate}
\item D. Negative cosine.
\item A. Coordinates by projection.
\item C. $MI^2=\dfrac{AB^2}{4}$.
\item B. Bilinearity: $2(-1)3$.
\item D. $BC^2=8$.
\end{enumerate}
\subsection*{True or false}
\begin{enumerate}
\item False: $u\cdot(-u)=-\norm u^2<0$ for $u\ne0$.
\item True: its product with every vector is 0.
\item False: $\pt{1;0}$ and $\pt{0;1}$ are nonzero and orthogonal.
\end{enumerate}
\paragraph{Differentiated teaching.} For a technical block, offer only the first step (domain, factor or reference sample space), then withdraw the help. For an independent student, ask for a second method and the treatment of a limiting case. Orally, allow an initial reformulation in the language the student has mastered, then require a precise mathematical sentence in French.
